Este capítulo revisa los antecedentes tecnológicos relevantes para el desarrollo del proyecto: las bases de datos de pictogramas disponibles, las técnicas existentes para generar texto natural a partir de palabras, las aplicaciones actuales de comunicación aumentativa y los motores de síntesis de voz.

\section{Bases de datos de pictogramas}

\textit{TODO: Esta sección requiere investigación adicional para comparar las bases de datos de pictogramas disponibles (ARASAAC, Mulberry, Sclera, etc.) y justificar la elección de ARASAAC.}

Desde la perspectiva de los recursos, varios trabajos han abordado la integración de pictogramas con recursos lingüísticos computacionales. Schwab et al.~\cite{schwab2020arasaac} propusieron vincular los pictogramas de ARASAAC con WordNet, creando una base de datos que relaciona símbolos con definiciones semánticas estructuradas. Este tipo de integración es fundamental para sistemas que necesitan razonar sobre el significado de los pictogramas, aunque los autores identificaron desafíos como la polisemia y la dependencia contextual de ciertas palabras.

\section{Técnicas de generación de texto natural}

El problema central que aborda este proyecto ---transformar secuencias de pictogramas en oraciones naturales--- ha sido estudiado desde múltiples perspectivas en la literatura. Este proceso involucra decisiones complejas: conjugación verbal, concordancia, inserción de artículos y preposiciones, e incluso reformulación semántica para producir expresiones más naturales.

Es importante señalar que los pictogramas en sistemas SAAC llevan asociada una etiqueta textual (por ejemplo, el pictograma de una casa tiene asociada la palabra ``casa''). Por tanto, el problema de generación \textit{picto-to-text} se reduce, en la práctica, a generar oraciones naturales a partir de una secuencia de palabras. Esta observación permite ampliar el alcance de la revisión a técnicas de generación de lenguaje natural que, aunque no hayan sido diseñadas específicamente para pictogramas, resuelven el mismo problema subyacente.

La literatura distingue dos tareas relacionadas pero distintas. Por un lado, la traducción de texto a pictogramas (\textit{text-to-picto}), que permite a una persona sin dificultades comunicativas expresarse en un formato comprensible para usuarios de SAAC. Por otro, la generación de texto a partir de pictogramas (\textit{picto-to-text}), que es el foco de este trabajo. La iniciativa ImageCLEF ToPicto~\cite{imageclef2024topicto}, por ejemplo, propuso en 2024 un desafío específico para la traducción de texto y voz a secuencias de pictogramas en francés, evidenciando el interés creciente en este campo.

\subsection{Técnicas basadas en gramáticas}

La solución más simple consiste en concatenar las etiquetas de los pictogramas seleccionados, sin procesamiento lingüístico alguno. Esta aproximación tiene ventajas evidentes: es predecible, fácil de implementar y no introduce errores de procesamiento. Sin embargo, produce secuencias como ``yo querer ir parque'' que, aunque comprensibles, imponen al interlocutor la tarea de reconstruir mentalmente la oración correcta.

Un paso intermedio lo constituyen los sistemas basados en plantillas. Estas aproximaciones definen estructuras predefinidas (por ejemplo, \texttt{[SUJETO] quiere [OBJETO]}) y rellenan los huecos con las palabras correspondientes. Las plantillas garantizan corrección gramatical dentro de su ámbito, pero adolecen de rigidez: cubren solo las estructuras anticipadas por los diseñadores, y cualquier combinación no prevista produce resultados pobres o nulos.

Una aproximación más sofisticada consiste en definir gramáticas formales que capturen las reglas de construcción de oraciones. El sistema puede entonces generar texto aplicando estas reglas a las palabras de entrada, tratando el problema como una tarea de realización superficial (\textit{surface realization}) en el campo de la Generación de Lenguaje Natural (GLN). La biblioteca SimpleNLG~\cite{gatt2009simplenlg}, originalmente desarrollada para inglés, ha sido adaptada a múltiples idiomas incluyendo español~\cite{ramos2017simplenlg, garciamendez2019spanish}. Para el ámbito de la comunicación aumentativa, García-Méndez et al.~\cite{garciamendez2018elsa} desarrollaron ELSA, el primer léxico español con información morfológica, sintáctica y semántica derivada de pictogramas ARASAAC. Combinado con SimpleNLG-ES, este recurso permite expandir selecciones de pictogramas a oraciones correctas, gestionando automáticamente la inserción de artículos, la selección de preposiciones y la conjugación verbal. El enfoque gramatical produce resultados predecibles y permite control fino sobre la generación, aunque su cobertura se limita a las estructuras previstas por las reglas.

Estas gramáticas operan en el plano sintáctico: garantizan que una oración esté bien formada, pero no verifican si tiene sentido. Para abordar esta limitación existen las restricciones de selección (\textit{selectional restrictions}), que asocian a cada predicado condiciones sobre el tipo semántico de sus argumentos~\cite{resnik1996selectional}. Por ejemplo, el verbo ``comer'' requiere un sujeto animado y un objeto comestible, lo que permite aceptar ``el niño come pan'' pero rechazar ``la mesa come pan''. La combinación de una gramática formal con restricciones de selección produce un sistema que verifica tanto la corrección sintáctica como la plausibilidad semántica de las oraciones generadas.

Algunos trabajos han aplicado estos principios al diseño de gramáticas específicas para comunicación aumentativa. Martínez-Santiago et al.~\cite{martinez2016pictogrammar} propusieron Pictogrammar, un sistema centrado en español que combina una ontología formal con la biblioteca Grammatical Framework para la realización morfosintáctica, restringiendo la salida a oraciones sintáctica y semánticamente válidas. Pereira et al.~\cite{pereira2020semanticgrammar} desarrollaron una gramática semántica para SAAC basada en \textit{Colourful Semantics} ---un marco terapéutico que codifica roles temáticos mediante colores (naranja para el agente, amarillo para la acción, verde para el objeto)--- combinada con análisis de dependencias y etiquetado de roles semánticos. Su evaluación sobre más de 1.200 oraciones sujeto-verbo-objeto alcanzó una precisión media del 90\% en la reconstrucción de oraciones a partir de pictogramas.

\subsection{Técnicas basadas en aprendizaje automático}

\textit{TODO: Esta sección se desarrollará tras la investigación de técnicas basadas en embeddings y modelos de lenguaje (LLMs) para generación de oraciones naturales a partir de palabras clave.}

\section{Aplicaciones actuales de comunicación aumentativa}

\textit{TODO: Esta sección revisará las aplicaciones SAAC disponibles actualmente para identificar sus deficiencias y oportunidades de mejora.}

Aplicaciones comerciales como \textit{LetMeTalk}~\cite{letmetalk2024} o \textit{Proloquo2Go}~\cite{proloquo2go2024} ofrecen tableros de pictogramas con síntesis de voz, permitiendo al usuario seleccionar símbolos que el sistema convierte en audio. Sin embargo, se limitan a concatenar las etiquetas sin procesamiento lingüístico.

El sistema Araboard~\cite{baldassarri2014aac}, desarrollado en la Universidad de Zaragoza, explora el enfoque de plantillas para el español en el contexto de comunicación aumentativa.

García-Méndez et al.~\cite{garciamendez2018aac} propusieron un sistema basado en SimpleNLG orientado específicamente a comunicación aumentativa, capaz de inferir automáticamente preposiciones y otros elementos a partir de información semántica, alcanzando una precisión del 77\% en la generación correcta de oraciones. El proyecto COMPANSION~\cite{demasco1992compansion}, desarrollado en los años 90, ya exploraba la expansión de mensajes comprimidos a oraciones completas mediante gramáticas de unificación, aunque su dependencia de marcadores léxicos explícitos y su elevado coste computacional limitaron su aplicabilidad práctica.

Hervás et al.~\cite{hervas2020predictivecomposition} presentan \textit{PictoEditor}, una aplicación en tablet que incorpora predicción en tableros pictográficos. En lugar de intentar ``traducir'' la frase completa, el sistema optimiza el flujo de selección de pictogramas anticipando qué símbolos son más probables a continuación. Sus resultados sugieren que la predicción basada en frecuencia mejora la disponibilidad inmediata del pictograma deseado, reduciendo la fricción de interacción~\cite{curtis2022stateofart}.

\section{Motores de síntesis de voz}

La generación de texto es solo una parte del problema; convertir ese texto en voz natural es igualmente importante para que el usuario pueda comunicarse de forma efectiva con su entorno.

Los primeros sistemas AAC utilizaban voces sintetizadas mediante concatenación de fonemas, produciendo un habla robótica e impersonal. Los usuarios frecuentemente reportaban que estas voces no les representaban, generando resistencia a su uso.

Los avances recientes en síntesis neuronal de voz han transformado este panorama. Motores como los ofrecidos por Google, Amazon u OpenAI producen locuciones que, en muchos contextos, resultan prácticamente indistinguibles del habla humana. Además, permiten personalización de parámetros como tono, velocidad y estilo, ofreciendo a los usuarios mayor capacidad de expresión~\cite{xu2025speakease}.

\textit{TODO: Comparativa de motores TTS disponibles (Google Cloud TTS, Amazon Polly, OpenAI TTS, etc.) para seleccionar el más adecuado para el proyecto.}
